% Отчёт по лабораторной работе №2.
\reportelement{Введение}

Цель лабораторной работы — освоить создание собственного Docker-образа
приложения, оптимизацию сборки, исключение лишних файлов из контекста
сборки, настройку проверки работоспособности контейнера и тегирование
готового образа.

В работе использовались Git, Docker Engine 28.5.2 и Debian GNU/Linux 13
в среде WSL2. Приложение реализовано на FastAPI и запускается через Uvicorn.

\begin{table}[H]
  \caption{Персональные параметры}\label{tab:parameters}
  \centering
  \begin{tabularx}{\textwidth}{|p{0.46\textwidth}|X|}
    \hline
    Параметр & Значение\\ \hline
    ФИО & Аннин Макар Дмитриевич\\ \hline
    Номер ИСУ & 465018\\ \hline
    Персональный порт & 5448\\ \hline
    Персональный префикс & \texttt{annin}\\ \hline
    Секретный ключ & \texttt{secret18annin}\\ \hline
    Группа & К3320\\ \hline
    Операционная система & Debian GNU/Linux 13 (WSL2)\\ \hline
    Архитектура & x86\_64\\ \hline
    Docker Engine & 28.5.2\\ \hline
    Git & 2.47.3\\ \hline
  \end{tabularx}
\end{table}

Исходный код и конфигурация работы размещены в публичном форке курса,
в ветке \texttt{lab-2}:
\url{https://github.com/VAPPROSS2/course-cont-orch-app/tree/lab-2}.

\chapter{Ход выполнения}

\section{Этап 1. Первый рабочий образ}

Для приложения был подготовлен Dockerfile на основе
\texttt{python:3.11-slim}. В нём заданы метаданные автора, рабочая
директория \texttt{/app}, переменная окружения \texttt{PORT=5448},
установка зависимостей и команда запуска Uvicorn.

Образ был собран командой:

\begin{code}
docker build -t annin-assistant:v1 .
\end{code}

Сборка завершилась успешно. В результате создан образ
\texttt{annin-assistant:v1}. Время первой сборки составило около
58,9 секунды. При сборке Docker вывел предупреждение
\texttt{JSONArgsRecommended}, относящееся к shell-форме команды CMD;
оно не помешало запуску приложения.

\reportscreenshot{images/step-01-build.png}{Успешная сборка первого Docker-образа}{fig:first-build}

Контейнер был запущен с публикацией порта 5448 и персональными
переменными окружения:

\begin{code}
docker run -d --name annin-assistant-lab2 \
  -p 5448:5448 \
  -e APP_OWNER=annin \
  -e APP_SECRET="$APP_SECRET" \
  annin-assistant:v1
\end{code}

Проверка главной страницы вернула приветствие для владельца
\texttt{annin}. Запрос \texttt{/me} вернул владельца и префикс секрета.
Запрос к \texttt{/secret/secret18annin} вернул флаг
\texttt{LAB2\_SUCCESS\_ANNIN}. Для проверки заметок была создана запись,
после чего GET-запрос вернул список заметок с количеством записей.

\reportscreenshot{images/step-01-root-secret.png}{Проверка главной страницы и секретного эндпоинта}{fig:root-secret}
\reportscreenshot{images/step-01-me.png}{Проверка персональных переменных окружения}{fig:me}
\reportscreenshot{images/step-01-notes.png}{Создание и получение заметок}{fig:notes}

\section{Этап 2. Оптимизация Docker-образа}

В файл \texttt{Dockerfile.optimized} добавлено отдельное копирование
\texttt{requirements.txt} до копирования исходного кода. Благодаря этому
слой установки зависимостей может повторно использоваться из кеша,
если изменяется только код приложения. Для установки пакетов применён
параметр \texttt{--no-cache-dir}, исключающий сохранение кеша pip в образе.

После сборки были получены следующие размеры:

\begin{table}[H]
  \caption{Сравнение размеров образов}\label{tab:image-sizes}
  \centering
  \begin{tabularx}{\textwidth}{|X|r|}
    \hline
    Образ & Размер\\ \hline
    \texttt{annin-assistant:v1} & 188 MB\\ \hline
    \texttt{annin-assistant:v2} & 175 MB\\ \hline
  \end{tabularx}
\end{table}

Оптимизированный образ оказался меньше на 13 MB, что составляет
примерно 6,9\% от первоначального размера.

\reportscreenshot{images/step-02-compare.png}{Сравнение размеров образов v1 и v2}{fig:compare}

\section{Этап 3. Исключение лишних файлов с помощью .dockerignore}

Создан файл \texttt{.dockerignore}, исключающий кеш Python, каталоги
виртуальных окружений, Git-метаданные, файл \texttt{.env}, логи,
Dockerfile и README из контекста сборки.

После добавления файла образ \texttt{annin-assistant:v3} был собран
повторно. Его размер составил 188 MB. В данном проекте размер совпал
с v1, поскольку Dockerfile копирует только \texttt{requirements.txt}
и каталог \texttt{app/}; исключённые файлы и ранее не попадали в
финальные слои образа.

Наличие каталога \texttt{.git} в контексте сборки нежелательно:
история репозитория может содержать удалённые файлы, конфигурации
или случайно сохранённые секреты. Исключение \texttt{.git} снижает
риск попадания такой информации в образ.

\reportscreenshot{images/step-03-dockerignore.png}{Образы v1, v2 и v3 после добавления .dockerignore}{fig:dockerignore}

\section{Этап 4. Настройка Healthcheck}

В \texttt{Dockerfile.optimized} добавлена установка утилиты curl
и инструкция HEALTHCHECK. Проверка выполняется обращением к
\texttt{http://localhost:\${PORT}/health} с заданными интервалом,
таймаутом, периодом запуска и числом повторных попыток.

Финальный образ собран командой:

\begin{code}
docker build -t annin-assistant:final -f Dockerfile.optimized .
\end{code}

После запуска контейнера \texttt{annin-assistant-final} и ожидания
инициализации команда \texttt{docker ps} показала статус
\texttt{Up (healthy)}. Это подтверждает, что Docker успешно выполняет
проверку эндпоинта здоровья приложения.

\reportscreenshot{images/step-04-health.png}{Статус финального контейнера healthy}{fig:health}

\section{Этап 5. Тегирование и финальный артефакт}

Финальному образу назначены дополнительные теги версии и даты сборки:

\begin{code}
docker tag annin-assistant:final annin-assistant:1.0.0
docker tag annin-assistant:final annin-assistant:$(date +%Y%m%d)
\end{code}

В результате созданы теги \texttt{final}, \texttt{1.0.0} и
\texttt{20261118}. Дополнительно сохранены образы \texttt{v1},
\texttt{v2} и \texttt{v3}.

Финальная команда вывела дату и время сборки, имя пользователя и
хоста, а также список тегов с размерами образов. В окружении выполнения
получено значение \texttt{Makar@HOME-PC}. Размер финального образа
с установленным curl составил 189 MB.

\reportscreenshot{images/step-05-artifact.png}{Финальный артефакт: дата, пользователь, хост, теги и размеры}{fig:artifact}

\reportelement{Заключение}

В ходе лабораторной работы был создан и запущен персонализированный
Docker-образ FastAPI-приложения. Проверены основные HTTP-эндпоинты,
переменные окружения и секретный эндпоинт. Оптимизация установки
зависимостей уменьшила размер образа с 188 MB до 175 MB. Добавлен
\texttt{.dockerignore}, настроен Healthcheck, подтверждён статус
контейнера \texttt{healthy}. Финальному образу назначены семантический
тег и тег с датой сборки.

Поставленная цель достигнута: процесс сборки, оптимизации, проверки
работоспособности и тегирования Docker-образа выполнен.

\reportelement{Список использованных источников}

\begin{enumerate}
  \item Docker Documentation. Dockerfile reference.
  URL: \url{https://docs.docker.com/reference/dockerfile/}
  (дата обращения: 18.11.2026).
  \item Docker Documentation. Build context and .dockerignore.
  URL: \url{https://docs.docker.com/build/concepts/context/}
  (дата обращения: 18.11.2026).
  \item Учебный репозиторий курса «Контейнеризация и оркестрация приложений».
  Лабораторная работа №2.
  URL: \url{https://github.com/itmo-faculty-applied-informatics/course-cont-orch-app/blob/main/labs/02-docker-image.md}
  (дата обращения: 18.11.2026).
